\clearpage\begingroup\NoSourceHead\SourceLeaf{521b}
\renewcommand{\thefootnote}{\arabic{footnote})}\setcounter{footnote}{19}
\ArithmeticHead{64}{L}{BOLZA, GAUSS UND DIE VARIATIONSRECHNUNG.}
\noindent\hspace*{\parindent}13. Diese Sätze von \textsc{Hamilton} sind mit den oben angeführten Sätzen
von \textsc{Gauss}, \textsc{Beltrami}, \textsc{Hilbert} und \textsc{Kneser} zusammen als Spezialfälle in den
folgenden allgemeinen Sätzen\footnote{Für $n=3$ bereits bei \textsc{Hamilton}, loc. cit. 18); wegen der Verallgemeinerung der Theorie auf das
\textsc{Lagrange}sche Problem vgl.: A. \textsc{Mayer}, \textit{Ueber den Hilbertschen Unabhängigkeitssatz in der Theorie der
Maxima und Minima der einfachen Integrale}, Leipziger Berichte, Bd. 55 (1903), S. 131—145 und Bd. 57
(1905), S. 49—67, 313—314; D. \textsc{Hilbert}, \textit{Zur Variationsrechnung}, Göttinger Nachrichten 1905, S. 159—180;
H. \textsc{Hahn}, \textit{Ueber den Zusammenhang zwischen den Theorien der zweiten Variation und der Weierstrassschen
Theorie der Variationsrechnung}, Rendiconti del Circolo Matematico di Palermo, Bd. 29 (1910), S. 49—78;
O. \textsc{Bolza}, \textit{Ueber den Hilbertschen Unabhängigkeitssatz beim Lagrangeschen Variationsproblem}, Rendiconti
del Circolo Matematico di Palermo, Bd. 31 (1911), S. 257—272; J. \textsc{Radon}, \textit{Zur Theorie der Mayerschen
Felder beim Lagrangeschen Variationsproblem}, Wiener Berichte, Bd. 120 (1911), S. 1337—1360. Vergl. auch
die Darstellung im ersten Teil der Dissertation von G. \textsc{Prange}, \textit{Die Hamilton-Jacobische Theorie für Doppelintegrale
(mit einer Uebersicht der Theorie für einfache Integrale)}, Göttingen 1915.} über die Extremalen des Integrals
\clearpage\SourceLeaf{522}
\ArithmeticHead{65}{R}{BEZIEHUNGEN ZUM UNABHÄNGIGKEITSSATZ.}
\begin{equation*}\tag{19}
\int F(x_1,x_2,\ldots,x_n,x'_1,x'_2,\ldots,x'_n)\,dt\equiv\int F(x,x')\,dt
\end{equation*}
enthalten: Damit eine $(n-1)$-parametrige Extremalenschar des Integrals (19),
die ein Feld bildet, eine Transversalhyperfläche besitzt, muss der Differentialausdruck
\begin{equation*}\tag{34}
\sum_{i=1}^n F_{x'_i}(x,p)\,dx_i,
\end{equation*}
in dem $p_1,p_2,\ldots,p_n$ die »Richtungskosinus« der durch den Punkt $x_1,x_2,\ldots,x_n$
gehenden Feldextremale sind, das vollständige Differential einer Funktion
$V(x_1,x_2,\ldots,x_n)$ sein, sodass also
\begin{equation*}\tag{35}
F_{x'_k}(x,p)=\frac{\partial V}{\partial x_k},\qquad k=1,2,\ldots,n.
\end{equation*}

Ist diese Bedingung erfüllt, in welchem Fall das Feld ein \textsc{Mayer}sches genannt
wird, so schneiden die sämtlichen Hyperflächen $V(x_1,x_2,\ldots,x_n)=\text{const.}$
die Extremalenschar transversal. Die Funktion $V$ ist mit dem Feldintegral
gerechnet von der Hyperfläche $V=0$ identisch. Die Richtungskosinus des
Feldes $p_1,p_2,\ldots,p_n$ genügen den die Gleichungen (25) und (33) als Spezialfälle
enthaltenden partiellen Differentialgleichungen
\begin{equation*}\tag{36}
\sum_{i=1}^n p_i\left(\frac{\partial F_{x'_k}(x,p)}{\partial x_i}-\frac{\partial F_{x'_i}(x,p)}{\partial x_k}\right)=0.
\end{equation*}

Die den partiellen Differentialgleichungen (26), (29), (32) entsprechende
partielle Differentialgleichung ergibt sich durch Elimination von $p_1,p_2,\ldots,p_n$
aus den $n$ Gleichungen (35) zusammen mit der Relation: $p_1^2+p_2^2+\cdots+p_n^2=1$.

Die angeführten Sätze sind als Spezialfälle in entsprechenden allgemeineren
über das sogenannte \textsc{Lagrange}sche Problem enthalten, auf das sich die in
Fussnote 20) zitierten Arbeiten beziehen. Eine noch weitergehende Verallgemeinerung
auf das sogenannte \textsc{Mayer}sche Problem hat neuerdings \textsc{Kneser}\begingroup\renewcommand{\thefootnote}{20a)}\footnote{A. \textsc{Kneser}, \textit{Beiträge zur Theorie der Variationsrechnung: Die Methode von Weierstrass im
Zusammenhang mit der Jacobi-Hamiltonschen und einer Integrationstheorie von Cauchy}, Archiv der Mathematik
und Physik (3), Bd. 24 (1915), S. 26—57. Nach G. \textsc{Prange} (Habilitationsrede, loc. cit. 19), S. 72)
hat sogar schon \textsc{Hamilton} seine allgemeine Methode auf ein \textsc{Mayer}sches Variationsproblem angewandt in
der Arbeit \textit{Calculus of principal relations}, Reports of the British Association for the advanc. of science 5
(1836, pt. 2), S. 41—44.}\endgroup\ gegeben.
\par\smallskip\noindent{\small X 2 Abh. 5.}\hfill{\small 9}\par
\clearpage\endgroup
